% ============================================================
%  sections/03_dataset.tex
%  3. Xây dựng Dataset  (~0.75 trang)
% ============================================================

\section{Xây Dựng Bộ Dữ Liệu}
\label{sec:dataset}

\subsection{Dữ Liệu Nguồn (Văn Bản Con Người)}

Để phản ánh thực tế đa dạng của mạng xã hội, chúng tôi chọn bốn corpus
MXH thực tế, mỗi corpus đại diện cho một ngôn ngữ và nền tảng khác nhau
(Bảng~\ref{tab:sources}).

\begin{table}[h]
\centering
\caption{Corpus dữ liệu nguồn (văn bản con người)}
\label{tab:sources}
\begin{tabular}{llll}
\toprule
\textbf{Ngôn ngữ} & \textbf{Corpus} & \textbf{Nền tảng} & \textbf{Số bài} \\
\midrule
EN & Sentiment140    & Twitter & 2.000 \\
VI & VSmec           & Facebook/MXH VN & 2.000 \\
ZH & Weibo Senti 100K & Weibo  & 2.000 \\
AR & AJGT            & Twitter & 2.000 \\
\bottomrule
\end{tabular}
\end{table}

Từ mỗi corpus, chúng tôi lấy mẫu ngẫu nhiên có phân tầng 2.000 bài
đăng, đảm bảo đa dạng về chủ đề và phong cách. Các bài đăng trùng lặp
và rỗng được loại bỏ trước khi tạo văn bản máy.

\subsection{Tạo Văn Bản Máy (Machine Text)}

Mỗi bài đăng con người được paraphrase 1:1 bởi một LLM phù hợp với
ngôn ngữ đó thông qua API (Bảng~\ref{tab:llms}). Chúng tôi sử dụng
prompt dạng paraphrase thay vì generation thuần túy nhằm giữ nguyên
chủ đề và ngữ cảnh, đồng thời tạo ra cặp (human, machine) có nội dung
tương đồng --- kịch bản thực tế nhất của MGT trên MXH.

\begin{table}[h]
\centering
\caption{LLM được sử dụng để tạo văn bản máy}
\label{tab:llms}
\begin{tabular}{llll}
\toprule
\textbf{Ngôn ngữ} & \textbf{Mô hình} & \textbf{Nhà cung cấp API} \\
\midrule
EN & LLaMA 3.1-8B-Instant      & Groq API \\
VI & GPT-OSS 20B               & NVIDIA NIM \\
ZH & Qwen3-80B-A3B-Instruct    & NVIDIA NIM \\
AR & LLaMA 4 Maverick 17B      & NVIDIA NIM \\
\bottomrule
\end{tabular}
\end{table}

Các bài đăng không được xử lý thành công (lỗi API, phản hồi rỗng hoặc
quá ngắn) bị loại bỏ. Đối với tiếng Việt, 7 bài bị loại do lỗi API
và vi phạm chính sách nội dung, dẫn đến 3.986 bài thay vì 4.000.

\subsection{Thống Kê Bộ Dữ Liệu Cuối}

Bộ dữ liệu cuối gồm \textbf{15.986 bài đăng} (7.993 human / 7.993
machine), được phân chia 80/20 theo phân tầng ngôn ngữ và nhãn
(Bảng~\ref{tab:dataset_stats}).

\begin{table}[h]
\centering
\caption{Thống kê bộ dữ liệu MultiSocial-4L}
\label{tab:dataset_stats}
\begin{tabular}{lrrrr}
\toprule
\textbf{Ngôn ngữ} & \textbf{Tổng} & \textbf{Train} & \textbf{Test} & \textbf{Ghi chú} \\
\midrule
EN & 4.000 & 3.200 & 800 & \\
VI & 3.986 & 3.189 & 797 & 7 bài bị loại \\
ZH & 4.000 & 3.200 & 800 & \\
AR & 4.000 & 3.200 & 800 & \\
\midrule
\textbf{Tổng} & \textbf{15.986} & \textbf{12.789} & \textbf{3.197} & \\
\bottomrule
\end{tabular}
\end{table}

Mỗi bản ghi gồm các trường: \texttt{text}, \texttt{label} (0=human,
1=machine), \texttt{language}, \texttt{length}, \texttt{source} và
\texttt{split}. Dữ liệu được cân bằng hoàn toàn (50\% human / 50\%
machine) theo từng ngôn ngữ.
